\chapter{Polygons and Area}

\chapterintro{Polygons are the shapes a renderer actually fills, and area is the quantity that tells you how much it costs to fill them. The shoelace formula in this chapter computes any polygon's area from its vertices in one loop, and its sign tells you which way the vertices wind, which is how a graphics pipeline decides what is a front face.}

\section{Core Ideas}

\paragraph{Angle sums} The interior angles of an $n$-sided polygon sum to $(n - 2) \cdot 180\degs$, because it splits into $n - 2$ triangles from one vertex. Each interior angle of a regular $n$-gon is that sum divided by $n$. The exterior angles of any convex polygon sum to $360\degs$, one full turn.

\paragraph{Quadrilaterals} A parallelogram has both pairs of opposite sides parallel; its opposite sides are equal and its diagonals bisect each other. A rectangle is a parallelogram with right angles; a rhombus is one with equal sides; a square is both. A trapezoid has exactly one pair of parallel sides. A kite has two pairs of adjacent equal sides and perpendicular diagonals.

\paragraph{Area formulas} Rectangle $lw$. Parallelogram $bh$, with $h$ perpendicular to the base. Triangle $\tfrac{1}{2}bh$. Trapezoid $\tfrac{1}{2}(b_1 + b_2)h$. Rhombus or kite $\tfrac{1}{2}d_1 d_2$ from the diagonals. Regular polygon $\tfrac{1}{2}aP$, where $a$ is the apothem (center to the middle of a side) and $P$ the perimeter.

\paragraph{The shoelace formula} For a simple polygon with vertices $(x_1, y_1), \ldots, (x_n, y_n)$ in order, with $(x_{n+1}, y_{n+1})$ meaning the first vertex again, the signed area is given below. It is positive when the vertices run counterclockwise and negative when clockwise; take the absolute value for the geometric area.

\begin{keybox}{Shoelace formula}
\[
A = \frac{1}{2}\sum_{i=1}^{n}\left(x_i\, y_{i+1} - x_{i+1}\, y_i\right)
\]
\end{keybox}

\paragraph{Scaling} Scaling a figure by $k$ multiplies lengths and perimeters by $k$ and areas by $k^2$. Doubling a sprite's size quadruples the pixels drawn.

\paragraph{Bounding boxes} The axis-aligned bounding box of a polygon runs from its smallest to largest $x$ and $y$. Two shapes whose boxes do not overlap cannot touch, which is why collision systems test boxes first and polygons only afterward.

\section{Worked Examples}

\begin{example}
Find the sum of the interior angles of a hexagon and each angle of a regular hexagon.
\solution $(6 - 2) \cdot 180\degs = 720\degs$; each angle of a regular hexagon is $720\degs / 6 = 120\degs$.
\end{example}

\begin{example}
Find the area of a trapezoid with bases $8$ and $12$ and height $5$.
\solution $\tfrac{1}{2}(8 + 12)(5) = 50$.
\end{example}

\begin{example}
Use the shoelace formula on the triangle $(0, 0)$, $(4, 0)$, $(0, 3)$.
\solution The terms are $0 \cdot 0 - 4 \cdot 0 = 0$, $4 \cdot 3 - 0 \cdot 0 = 12$, and $0 \cdot 0 - 0 \cdot 3 = 0$. Half the sum is $6$, and the positive sign says the vertices run counterclockwise.
\end{example}

\section{Exercises}

\begin{multicols}{2}
\begin{exercises}
\item Find the interior angle sum of an octagon, each interior angle of a regular octagon, and each exterior angle.
\item A regular polygon has interior angles of $150\degs$. How many sides does it have?
\item Find the area of a parallelogram with base $9$ and height $4$, and of a triangle with base $10$ and height $7$.
\item Find the area of a trapezoid with bases $6$ and $14$ and height $5$.
\item A rhombus has diagonals $10$ and $24$. Find its area and its side length.
\item Find the area of a regular hexagon with side $6$, exactly and to one decimal.
\item Use the shoelace formula on $(1, 1)$, $(5, 1)$, $(6, 4)$, $(2, 5)$.
\item A rectangle is scaled by $3$. By what factors do its perimeter and area change?
\item Find the area of the parallelogram $(0, 0)$, $(5, 0)$, $(7, 3)$, $(2, 3)$ using base times height, then confirm with the shoelace formula.
\item Find the bounding box of the triangle $(-2, 4)$, $(5, -1)$, $(3, 6)$ and its area, then the triangle's area by shoelace. What fraction of the box does the triangle fill?
\item \textbf{(proof)} Prove that the exterior angles of a convex $n$-gon sum to $360\degs$, using the interior angle sum.
\end{exercises}
\end{multicols}

\section{Programming Exercises}

\begin{codeexercises}
\item Write \lstinline!signed_area(points)! implementing the shoelace formula. Then write the function \lstinline!is_counterclockwise(points)!, which reads the sign of that result. Test with Exercises~6.7, 6.9 and 6.10, and with each polygon's vertices reversed.
\item Write \lstinline!regular_polygon(n, side)! that returns the $n$ vertices on a circle (use \lstinline!math.cos! and \lstinline!math.sin! with the angle $2\pi i / n$; Chapter~9 explains the radians). Compute its area with C6.1 and compare with $\tfrac{1}{2}aP$ for $n = 6$, side $6$.
\item Write \lstinline!bounding_box(points)! and \lstinline!boxes_overlap(a, b)!. Generate random triangles and confirm that whenever the boxes do not overlap, no vertex of one triangle lies inside the other.
\item Write \lstinline!interior_angle_sum(n)! and \lstinline!sides_from_interior_angle(deg)!, raising when the angle does not correspond to a whole number of sides. Test with Exercises~6.1 and 6.2.
\item Scale a polygon by $k$ using \lstinline!dilate! from Chapter~5 for $k = 1, 2, 3, 4, 5$, compute each area with C6.1, and print the ratio to the original to confirm it is $k^2$.
\end{codeexercises}
